\section{Introduction}
\label{sec:introduction}

Knowledge graphs have become the default representation for semantic
information on the web and in enterprise systems. Their foundations are simple:
every fact is a triple $(s, p, o)$ that asserts a binary relationship
between a subject, a predicate, and an object~\cite{manola2004rdf,sparql2013}. The
model is uniform, transparent, and has an enormous accumulated tooling base. It
is also a poor fit for large classes of real data.

Two mismatches recur across domains. First, many facts are not binary. A drug
interaction record carries a drug, a target, a dosage, a route, a study population,
and a confidence interval, all in one statement. A financial transaction pairs
a buyer, a seller, an instrument, a price, a currency, a venue, and a
timestamp. Encoding such facts as triples requires \emph{reification}: an
auxiliary event node is minted, each attribute becomes its own triple pointing
at that node, and reconstructing the original fact costs one join per attribute.
Second, the same fact often holds in some contexts and not in others. A
biological interaction that is valid \emph{in vitro} may fail \emph{in vivo};
a corporate policy that applies in the European subsidiary may not apply in
North America. Context is not first-class in the triple model. It is added
either through separate named graphs~\cite{carroll2005named}, through further
reification, or through ad-hoc conventions in the predicate URI.

Both problems have well-known workarounds, and mature triple stores handle
them competently at scale. The question we pose in this paper is whether there
is anything to gain from a database whose primary organising principle is not
the triple but the \emph{context}. Concretely, we ask whether one can build a
storage engine that (i)~stores each fact as an atomic record, without
reification, and (ii)~indexes those records by the context poset in a way that
answers entity-centric queries via a direct index lookup while remaining verifiably
equivalent to a conventional RDF-style store on the queries that both engines
can answer.

The mathematical vocabulary that fits this question naturally is that of
presheaves on a poset of contexts. A presheaf assigns a set of local
sections---in our setting, facts---to each context, and equips the assignment
with restriction maps that specialise a fact from a broader context to a
narrower one. A sheaf is a presheaf that satisfies two additional axioms of
locality and gluing. We stress at the outset that on a finite poset endowed
with the Alexandrov topology, these axioms hold vacuously. The system we describe
is therefore a \emph{presheaf-based} database in the strict technical sense; the
sheaf terminology is retained because the operational vocabulary of restriction,
stalks, and global sections is what our implementation actually exposes, not
because we claim a non-trivial sheaf-theoretic theorem.

\subsection*{Contributions}
\label{sec:contributions}

The contributions of this paper are, in order:

\begin{enumerate}
    \item A concise formal model of a context-indexed semantic fact store,
          expressed in the language of presheaves on a finite poset
          (Section~\ref{sec:mathematical_model}). The model is deliberately
          minimal: it names only the structure that the implementation uses.

    \item A prototype implementation, the \emph{Semantic Fact Database}
          (SFDB), that realises the model in roughly $16\,500$ lines of
          Python (Section~\ref{sec:implementation}). Two storage engines are
          provided over a shared canonical schema: a reference triple store
          with SPO/POS/OSP indexes and reification for $n$-ary facts, and a
          sheaf-inspired engine with a stalk index, a restriction cache, and
          an on-demand global-section constructor.

    \item A canonical intermediate representation together with a
          bidirectional mapping to and from each engine's native format,
          used to establish semantic equivalence at every insert and every
          query (Section~\ref{sec:correctness}). Cross-engine equivalence is
          enforced by the test suite ($416$ passing tests) and re-checked
          on every benchmark run.

    \item A benchmark evaluation across five query classes (LOOKUP,
          GLOBAL, CONTEXT, bounded TEMPORAL, and unbounded TEMPORAL) and insert
          throughput, at \benchscalecount{} scales
          (\benchscales{} facts; Section~\ref{sec:evaluation}), with
          every query at every scale checked for three-way cross-engine
          result equivalence as part of the same run. The evaluation
          includes a \emph{storage-layer control}: a variant of the KG
          baseline with identical reification and query logic whose
          SQLite tables are replaced by in-process dict indexes, so that
          the comparison separates what context-indexed storage buys
          from what merely avoiding SQLite round trips buys
          (Section~\ref{sec:eval:control}), plus four external reference
          points: rdflib, a real independent Python library
          (Section~\ref{sec:eval:rdflib}), and three real production RDF
          stores accessed over SPARQL/HTTP, Apache Jena TDB2
          (Section~\ref{sec:eval:jena}), OpenLink Virtuoso
          (Section~\ref{sec:eval:virtuoso}), and Ontotext GraphDB
          (Section~\ref{sec:eval:graphdb}) --- with result counts
          verified against all four on every class at every scale. The
          sheaf-inspired engine is
          \rangelookupkg{} faster than the KG baseline on LOOKUP queries
          (\rangelookupkgmem{} against the control),
          \rangeglobalkg{} faster on GLOBAL scans
          (\rangeglobalkgmem{} against the control), and
          \rangecontextkg{} faster on CONTEXT queries --- the query shape
          the design is named for --- (\rangecontextkgmem{} against the
          control); against all three real production stores, SFDB remains
          faster on LOOKUP, GLOBAL, and CONTEXT at every scale but is overtaken on
          TEMPORAL from $10^4$ facts onward, a genuine
          crossover, reproduced independently against three unrelated
          optimisers, that we report as measured rather than as either a
          clean win or a hedge. The temporal-range
          advantages over the SQLite baseline do not survive the KG-mem control at
          all and are reported as a negative result, and insert
          throughput is slower than the write-through baseline at the
          smallest scale measured and faster at every larger one.
          The tradeoff we report is memory: SFDB's in-memory indexing
          structures use \rangememkg{} more resident memory per fact
          than the KG baseline's reified-triple storage. We also report
          that tightening the cross-engine verification harness during
          this evaluation caught a real semantic bug in the sheaf
          engine's LOOKUP path (Section~\ref{sec:correctness}), which we
          take as evidence for the verification methodology, not just
          for the numbers. Two further comparisons, a real-world case study
          (Section~\ref{sec:eval:wikidata}, \wikidatanumfacts{}
          genuine Wikidata facts) and the LUBM standard benchmark workload
          (Section~\ref{sec:eval:lubm}, up to \lubmlargescale{} facts),
          confirm
          the LOOKUP/GLOBAL/CONTEXT advantage on both while leaving TEMPORAL
          in a different place on each, and building the LUBM comparison
          surfaced a genuine CONTEXT-resolution performance bug
          (Section~\ref{sec:discussion:threats}) that had been present
          throughout the entire evaluation, requiring every CONTEXT number
          in this paper to be re-measured after the fix.

    \item A discussion of the design's realistic scope, including the
          honest observation that the sheaf condition is vacuous in our
          setting and therefore not the source of any of the performance
          differences we observe (Section~\ref{sec:discussion}). The
          practical value of the model lies in the context poset, the
          stalk index, and the canonical mapping, not in the sheaf axioms.
\end{enumerate}

\subsection*{What This Paper Is Not}

We do not claim that a sheaf-inspired database replaces production RDF
stores such as Apache Jena, GraphDB, or Virtuoso. Our KG baseline is a
research prototype implemented by the same authors as the SFDB engine, and
we call attention to this limitation throughout: the LOOKUP, GLOBAL, and
CONTEXT speedups against our own baselines
are best interpreted as measurements of what disciplined
context-indexed storage buys over a strict reification-plus-joins design at
comparable engineering quality, not as claims about the state of the art in
RDF engines. We do report direct comparisons against three real production
stores, Apache Jena TDB2 (Section~\ref{sec:eval:jena}), OpenLink
Virtuoso (Section~\ref{sec:eval:virtuoso}), and Ontotext GraphDB
(Section~\ref{sec:eval:graphdb}) --- SFDB is faster than all three
on LOOKUP, GLOBAL, and CONTEXT, and slower than all three on TEMPORAL from
$10^4$ facts onward --- and against LUBM
(Section~\ref{sec:eval:lubm}), the standard benchmark workload for RDF
stores, at scales up to \lubmlargescale{} facts, but not yet against
BSBM or WatDiv, or LUBM at the much larger scales published LUBM
results use; extending further in that direction remains future
work. We do not claim novelty in category theory. The formal machinery we use is
standard.

\subsection*{Paper Organisation}
\label{sec:organization}

Section~\ref{sec:related-work} places the work relative to existing knowledge
graph and category-theoretic database literature. Section~\ref{sec:background}
recalls the finite mathematical background we rely on. Section~\ref{sec:problem}
states the storage and equivalence problems formally.
Section~\ref{sec:mathematical_model} presents the formal model, together
with the definitions and lemmas that back it.
Section~\ref{sec:system-architecture} describes the system architecture, and
Section~\ref{sec:implementation} the implementation. Section~\ref{sec:evaluation}
reports the benchmark results, organised by query class.
Section~\ref{sec:discussion} interprets the results, and
Section~\ref{sec:limitations} states the limitations. Section~\ref{sec:future}
outlines future work, and Section~\ref{sec:conclusion} concludes.
